% Options for packages loaded elsewhere
\PassOptionsToPackage{unicode}{hyperref}
\PassOptionsToPackage{hyphens}{url}
%
\documentclass[
  ignorenonframetext,
]{beamer}
\usepackage{pgfpages}
\setbeamertemplate{caption}[numbered]
\setbeamertemplate{caption label separator}{: }
\setbeamercolor{caption name}{fg=normal text.fg}
\beamertemplatenavigationsymbolsempty
% Prevent slide breaks in the middle of a paragraph
\widowpenalties 1 10000
\raggedbottom
\setbeamertemplate{part page}{
  \centering
  \begin{beamercolorbox}[sep=16pt,center]{part title}
    \usebeamerfont{part title}\insertpart\par
  \end{beamercolorbox}
}
\setbeamertemplate{section page}{
  \centering
  \begin{beamercolorbox}[sep=12pt,center]{part title}
    \usebeamerfont{section title}\insertsection\par
  \end{beamercolorbox}
}
\setbeamertemplate{subsection page}{
  \centering
  \begin{beamercolorbox}[sep=8pt,center]{part title}
    \usebeamerfont{subsection title}\insertsubsection\par
  \end{beamercolorbox}
}
\AtBeginPart{
  \frame{\partpage}
}
\AtBeginSection{
  \ifbibliography
  \else
    \frame{\sectionpage}
  \fi
}
\AtBeginSubsection{
  \frame{\subsectionpage}
}
\usepackage{amsmath,amssymb}
\usepackage{lmodern}
\usepackage{iftex}
\ifPDFTeX
  \usepackage[T1]{fontenc}
  \usepackage[utf8]{inputenc}
  \usepackage{textcomp} % provide euro and other symbols
\else % if luatex or xetex
  \usepackage{unicode-math}
  \defaultfontfeatures{Scale=MatchLowercase}
  \defaultfontfeatures[\rmfamily]{Ligatures=TeX,Scale=1}
\fi
\usetheme[]{Montpellier}
% Use upquote if available, for straight quotes in verbatim environments
\IfFileExists{upquote.sty}{\usepackage{upquote}}{}
\IfFileExists{microtype.sty}{% use microtype if available
  \usepackage[]{microtype}
  \UseMicrotypeSet[protrusion]{basicmath} % disable protrusion for tt fonts
}{}
\makeatletter
\@ifundefined{KOMAClassName}{% if non-KOMA class
  \IfFileExists{parskip.sty}{%
    \usepackage{parskip}
  }{% else
    \setlength{\parindent}{0pt}
    \setlength{\parskip}{6pt plus 2pt minus 1pt}}
}{% if KOMA class
  \KOMAoptions{parskip=half}}
\makeatother
\usepackage{xcolor}
\IfFileExists{xurl.sty}{\usepackage{xurl}}{} % add URL line breaks if available
\IfFileExists{bookmark.sty}{\usepackage{bookmark}}{\usepackage{hyperref}}
\hypersetup{
  hidelinks,
  pdfcreator={LaTeX via pandoc}}
\urlstyle{same} % disable monospaced font for URLs
\newif\ifbibliography
\usepackage{color}
\usepackage{fancyvrb}
\newcommand{\VerbBar}{|}
\newcommand{\VERB}{\Verb[commandchars=\\\{\}]}
\DefineVerbatimEnvironment{Highlighting}{Verbatim}{commandchars=\\\{\}}
% Add ',fontsize=\small' for more characters per line
\usepackage{framed}
\definecolor{shadecolor}{RGB}{248,248,248}
\newenvironment{Shaded}{\begin{snugshade}}{\end{snugshade}}
\newcommand{\AlertTok}[1]{\textcolor[rgb]{0.94,0.16,0.16}{#1}}
\newcommand{\AnnotationTok}[1]{\textcolor[rgb]{0.56,0.35,0.01}{\textbf{\textit{#1}}}}
\newcommand{\AttributeTok}[1]{\textcolor[rgb]{0.77,0.63,0.00}{#1}}
\newcommand{\BaseNTok}[1]{\textcolor[rgb]{0.00,0.00,0.81}{#1}}
\newcommand{\BuiltInTok}[1]{#1}
\newcommand{\CharTok}[1]{\textcolor[rgb]{0.31,0.60,0.02}{#1}}
\newcommand{\CommentTok}[1]{\textcolor[rgb]{0.56,0.35,0.01}{\textit{#1}}}
\newcommand{\CommentVarTok}[1]{\textcolor[rgb]{0.56,0.35,0.01}{\textbf{\textit{#1}}}}
\newcommand{\ConstantTok}[1]{\textcolor[rgb]{0.00,0.00,0.00}{#1}}
\newcommand{\ControlFlowTok}[1]{\textcolor[rgb]{0.13,0.29,0.53}{\textbf{#1}}}
\newcommand{\DataTypeTok}[1]{\textcolor[rgb]{0.13,0.29,0.53}{#1}}
\newcommand{\DecValTok}[1]{\textcolor[rgb]{0.00,0.00,0.81}{#1}}
\newcommand{\DocumentationTok}[1]{\textcolor[rgb]{0.56,0.35,0.01}{\textbf{\textit{#1}}}}
\newcommand{\ErrorTok}[1]{\textcolor[rgb]{0.64,0.00,0.00}{\textbf{#1}}}
\newcommand{\ExtensionTok}[1]{#1}
\newcommand{\FloatTok}[1]{\textcolor[rgb]{0.00,0.00,0.81}{#1}}
\newcommand{\FunctionTok}[1]{\textcolor[rgb]{0.00,0.00,0.00}{#1}}
\newcommand{\ImportTok}[1]{#1}
\newcommand{\InformationTok}[1]{\textcolor[rgb]{0.56,0.35,0.01}{\textbf{\textit{#1}}}}
\newcommand{\KeywordTok}[1]{\textcolor[rgb]{0.13,0.29,0.53}{\textbf{#1}}}
\newcommand{\NormalTok}[1]{#1}
\newcommand{\OperatorTok}[1]{\textcolor[rgb]{0.81,0.36,0.00}{\textbf{#1}}}
\newcommand{\OtherTok}[1]{\textcolor[rgb]{0.56,0.35,0.01}{#1}}
\newcommand{\PreprocessorTok}[1]{\textcolor[rgb]{0.56,0.35,0.01}{\textit{#1}}}
\newcommand{\RegionMarkerTok}[1]{#1}
\newcommand{\SpecialCharTok}[1]{\textcolor[rgb]{0.00,0.00,0.00}{#1}}
\newcommand{\SpecialStringTok}[1]{\textcolor[rgb]{0.31,0.60,0.02}{#1}}
\newcommand{\StringTok}[1]{\textcolor[rgb]{0.31,0.60,0.02}{#1}}
\newcommand{\VariableTok}[1]{\textcolor[rgb]{0.00,0.00,0.00}{#1}}
\newcommand{\VerbatimStringTok}[1]{\textcolor[rgb]{0.31,0.60,0.02}{#1}}
\newcommand{\WarningTok}[1]{\textcolor[rgb]{0.56,0.35,0.01}{\textbf{\textit{#1}}}}
\setlength{\emergencystretch}{3em} % prevent overfull lines
\providecommand{\tightlist}{%
  \setlength{\itemsep}{0pt}\setlength{\parskip}{0pt}}
\setcounter{secnumdepth}{-\maxdimen} % remove section numbering
%\documentclass[unicode,10pt]{beamer}% 'unicode'が必要
\usepackage{luatexja}% 日本語を使う場合は必須
%\usepackage[japanese]{babel}	
%\usefonttheme[onlymath]{serif}
%\usepackage[T1]{fontenc} %これを使うと、ギリシャ文字が出ない
%\usepackage{textcomp}
%\usepackage[scale=1.0]{tgheros} % Sans serif font
%\usepackage[scaled]{beramono} % Monospace font
%\usepackage{luatexja-otf}
%\renewcommand{\kanjifamilydefault}{\gtdefault}
%\usepackage[haranoaji]{luatexja-preset}

% スライド番号の表示 %%%%%%%%%%%%%%%%%%%
\makeatletter
\setbeamertemplate{footline}{%
  \leavevmode%
  \hbox{%
  \begin{beamercolorbox}[wd=.5\paperwidth,ht=2.25ex,dp=1ex,right]{date in head/foot}%
    \insertframenumber{} \hspace*{2ex} % new version without total frames
  \end{beamercolorbox}}%
  \vskip0pt%
}
\makeatother
% スライド番号の表示 %%%%%%%%%%%%%%%%%%%
\ifLuaTeX
  \usepackage{selnolig}  % disable illegal ligatures
\fi

\title{第5章: 発見\\
(5.3. 空間データ)}
\subtitle{今井耕介 著\\
『社会科学のためのデータ分析入門 (QSS)』}
\author{}
\date{\vspace{-2.5em}2026-03-09}

\begin{document}
\frame{\titlepage}

\hypertarget{ux7a7aux9593ux30c7ux30fcux30bf}{%
\section{5.3 空間データ}\label{ux7a7aux9593ux30c7ux30fcux30bf}}

\begin{frame}{地理空間情報の分析}
\protect\hypertarget{ux5730ux7406ux7a7aux9593ux60c5ux5831ux306eux5206ux6790}{}
\begin{itemize}
\tightlist
\item
  \textbf{空間データ}: 位置情報（緯度・経度など）を持つデータ。
\item
  \textbf{目標}:
  現象の地理的な広がりを視覚化し、地域間のパターンや関係を理解する。
\item
  古典的な例: スノーによるコレラ発生源の特定。
\end{itemize}
\end{frame}

\hypertarget{rux306bux304aux3051ux308bux7a7aux9593ux30c7ux30fcux30bfux306eux6271ux3044}{%
\section{5.3.1
Rにおける空間データの扱い}\label{rux306bux304aux3051ux308bux7a7aux9593ux30c7ux30fcux30bfux306eux6271ux3044}}

\begin{frame}[fragile]{maps パッケージ: プロットのコード}
\protect\hypertarget{maps-ux30d1ux30c3ux30b1ux30fcux30b8-ux30d7ux30edux30c3ux30c8ux306eux30b3ux30fcux30c9}{}
\begin{itemize}
\tightlist
\item
  Rで地図を描画するための標準的なパッケージ。
\end{itemize}

\scriptsize

\begin{Shaded}
\begin{Highlighting}[]
\FunctionTok{library}\NormalTok{(maps)}
\CommentTok{\# 1. 米国地図の描画}
\FunctionTok{map}\NormalTok{(}\AttributeTok{database =} \StringTok{"usa"}\NormalTok{) }

\CommentTok{\# 2. 州都のデータ (us.cities) を読み込む}
\FunctionTok{data}\NormalTok{(us.cities) }
\CommentTok{\# capital == 2 が州都を表す}
\NormalTok{capitals }\OtherTok{\textless{}{-}} \FunctionTok{subset}\NormalTok{(us.cities, capital }\SpecialCharTok{==} \DecValTok{2}\NormalTok{) }

\CommentTok{\# 3. 州都の位置を地図上にプロット}
\CommentTok{\# cex は点の大きさ、pch は点の形}
\FunctionTok{points}\NormalTok{(}\AttributeTok{x =}\NormalTok{ capitals}\SpecialCharTok{$}\NormalTok{long, }\AttributeTok{y =}\NormalTok{ capitals}\SpecialCharTok{$}\NormalTok{lat,}
       \AttributeTok{cex =}\NormalTok{ capitals}\SpecialCharTok{$}\NormalTok{pop }\SpecialCharTok{/} \DecValTok{500000}\NormalTok{, }\AttributeTok{pch =} \DecValTok{19}\NormalTok{, }\AttributeTok{col =} \StringTok{"red"}\NormalTok{)}
\end{Highlighting}
\end{Shaded}

\normalsize
\end{frame}

\begin{frame}{maps パッケージ: 描画結果}
\protect\hypertarget{maps-ux30d1ux30c3ux30b1ux30fcux30b8-ux63cfux753bux7d50ux679c}{}
\scriptsize

\begin{center}\includegraphics[height=0.75\textheight]{QSS_Ch05_3_files/figure-beamer/unnamed-chunk-2-1} \end{center}
\normalsize
\end{frame}

\begin{frame}[fragile]{特定の地域の拡大: プロットのコード}
\protect\hypertarget{ux7279ux5b9aux306eux5730ux57dfux306eux62e1ux5927-ux30d7ux30edux30c3ux30c8ux306eux30b3ux30fcux30c9}{}
\begin{itemize}
\tightlist
\item
  \texttt{regions} 引数で特定の州などを指定できます。
\end{itemize}

\scriptsize

\begin{Shaded}
\begin{Highlighting}[]
\CommentTok{\# 1. カリフォルニア州のみを描画}
\FunctionTok{map}\NormalTok{(}\AttributeTok{database =} \StringTok{"state"}\NormalTok{, }\AttributeTok{regions =} \StringTok{"California"}\NormalTok{)}

\CommentTok{\# 2. カリフォルニア州の都市データを抽出}
\NormalTok{cal.cities }\OtherTok{\textless{}{-}} \FunctionTok{subset}\NormalTok{(us.cities, }\AttributeTok{subset =}\NormalTok{ (country.etc }\SpecialCharTok{==} \StringTok{"CA"}\NormalTok{))}

\CommentTok{\# 3. 人口上位7都市を特定し、プロット}
\NormalTok{top7 }\OtherTok{\textless{}{-}} \FunctionTok{order}\NormalTok{(cal.cities}\SpecialCharTok{$}\NormalTok{pop, }\AttributeTok{decreasing =} \ConstantTok{TRUE}\NormalTok{)[}\DecValTok{1}\SpecialCharTok{:}\DecValTok{7}\NormalTok{]}
\FunctionTok{points}\NormalTok{(}\AttributeTok{x =}\NormalTok{ cal.cities}\SpecialCharTok{$}\NormalTok{long[top7], }\AttributeTok{y =}\NormalTok{ cal.cities}\SpecialCharTok{$}\NormalTok{lat[top7], }\AttributeTok{pch =} \DecValTok{19}\NormalTok{)}

\CommentTok{\# 4. 都市名を追加}
\FunctionTok{text}\NormalTok{(}\AttributeTok{x =}\NormalTok{ cal.cities}\SpecialCharTok{$}\NormalTok{long[top7] }\SpecialCharTok{+} \DecValTok{2}\NormalTok{, }\AttributeTok{y =}\NormalTok{ cal.cities}\SpecialCharTok{$}\NormalTok{lat[top7],}
     \AttributeTok{label =}\NormalTok{ cal.cities}\SpecialCharTok{$}\NormalTok{name[top7], }\AttributeTok{cex =} \FloatTok{0.8}\NormalTok{)}
\end{Highlighting}
\end{Shaded}

\normalsize
\end{frame}

\begin{frame}{特定の地域の拡大: 描画結果}
\protect\hypertarget{ux7279ux5b9aux306eux5730ux57dfux306eux62e1ux5927-ux63cfux753bux7d50ux679c}{}
\scriptsize

\begin{center}\includegraphics[height=0.75\textheight]{QSS_Ch05_3_files/figure-beamer/unnamed-chunk-4-1} \end{center}
\normalsize
\end{frame}

\hypertarget{rux306bux304aux3051ux308bux8272ux306eux6307ux5b9a}{%
\section{5.3.2
Rにおける色の指定}\label{rux306bux304aux3051ux308bux8272ux306eux6307ux5b9a}}

\begin{frame}[fragile]{rgb() 関数と透明度}
\protect\hypertarget{rgb-ux95a2ux6570ux3068ux900fux660eux5ea6}{}
\begin{itemize}
\tightlist
\item
  \texttt{rgb()} で赤・緑・青の強度 (0〜1) と、透明度 (\texttt{alpha})
  を指定できます。
\end{itemize}

\scriptsize

\begin{Shaded}
\begin{Highlighting}[]
\CommentTok{\# 半透明の青 (alpha=0.5) を作成}
\NormalTok{blue.trans }\OtherTok{\textless{}{-}} \FunctionTok{rgb}\NormalTok{(}\AttributeTok{red =} \DecValTok{0}\NormalTok{, }\AttributeTok{green =} \DecValTok{0}\NormalTok{, }\AttributeTok{blue =} \DecValTok{1}\NormalTok{, }\AttributeTok{alpha =} \FloatTok{0.5}\NormalTok{)}

\CommentTok{\# 重なりを確認するためのプロット}
\FunctionTok{plot}\NormalTok{(}\AttributeTok{x =} \FunctionTok{c}\NormalTok{(}\DecValTok{1}\NormalTok{, }\FloatTok{1.2}\NormalTok{), }\AttributeTok{y =} \FunctionTok{c}\NormalTok{(}\DecValTok{1}\NormalTok{, }\DecValTok{1}\NormalTok{), }\AttributeTok{xlim =} \FunctionTok{c}\NormalTok{(}\FloatTok{0.5}\NormalTok{, }\DecValTok{2}\NormalTok{), }\AttributeTok{ylim =} \FunctionTok{c}\NormalTok{(}\FloatTok{0.5}\NormalTok{, }\FloatTok{1.5}\NormalTok{), }
     \AttributeTok{pch =} \DecValTok{16}\NormalTok{, }\AttributeTok{cex =} \DecValTok{10}\NormalTok{, }\AttributeTok{col =}\NormalTok{ blue.trans, }\AttributeTok{main =} \StringTok{"Transparency Example"}\NormalTok{)}
\end{Highlighting}
\end{Shaded}

\begin{center}\includegraphics[height=0.4\textheight]{QSS_Ch05_3_files/figure-beamer/unnamed-chunk-5-1} \end{center}
\normalsize
\end{frame}

\hypertarget{ux7c73ux56fdux5927ux7d71ux9818ux9078ux6319ux306eux30deux30c3ux30d7}{%
\section{5.3.3
米国大統領選挙のマップ}\label{ux7c73ux56fdux5927ux7d71ux9818ux9078ux6319ux306eux30deux30c3ux30d7}}

\begin{frame}[fragile]{pres08 データの読み込み (1)}
\protect\hypertarget{pres08-ux30c7ux30fcux30bfux306eux8aadux307fux8fbcux307f-1}{}
\begin{itemize}
\tightlist
\item
  民主党が勝った州と共和党が勝った州の名前を、それぞれ別々のリストにまとめます。
\end{itemize}

\scriptsize

\begin{Shaded}
\begin{Highlighting}[]
\CommentTok{\# 1. ローカルに保存したデータの読み込み（推奨）}
\NormalTok{pres08 }\OtherTok{\textless{}{-}} \FunctionTok{read.csv}\NormalTok{(}\StringTok{"pres08.csv"}\NormalTok{)}

\CommentTok{\# （参考）URLから直接読み込むことも可能}
\CommentTok{\# pres08 \textless{}{-} read.csv("https://ayumu{-}tanaka.github.io/QSS/QSS\_Data/pres08.csv")}
\end{Highlighting}
\end{Shaded}

\normalsize
\end{frame}

\begin{frame}[fragile]{州のリストを作成 (2)}
\protect\hypertarget{ux5ddeux306eux30eaux30b9ux30c8ux3092ux4f5cux6210-2}{}
\scriptsize

\begin{Shaded}
\begin{Highlighting}[]
\CommentTok{\# 1. オバマ（民主党）が勝った州の名前を取り出す}
\CommentTok{\# Obama \textgreater{} McCain という条件に合う state.name を抽出}
\NormalTok{blue\_states }\OtherTok{\textless{}{-}}\NormalTok{ pres08}\SpecialCharTok{$}\NormalTok{state.name[pres08}\SpecialCharTok{$}\NormalTok{Obama }\SpecialCharTok{\textgreater{}}\NormalTok{ pres08}\SpecialCharTok{$}\NormalTok{McCain]}

\CommentTok{\# 2. マケイン（共和党）が勝った州の名前を取り出す}
\CommentTok{\# McCain \textgreater{} Obama という条件に合う state.name を抽出}
\NormalTok{red\_states }\OtherTok{\textless{}{-}}\NormalTok{ pres08}\SpecialCharTok{$}\NormalTok{state.name[pres08}\SpecialCharTok{$}\NormalTok{McCain }\SpecialCharTok{\textgreater{}}\NormalTok{ pres08}\SpecialCharTok{$}\NormalTok{Obama]}

\CommentTok{\# リストの最初の方を確認}
\FunctionTok{head}\NormalTok{(blue\_states, }\AttributeTok{n =} \DecValTok{5}\NormalTok{)}
\end{Highlighting}
\end{Shaded}

\begin{verbatim}
## [1] "California"  "Colorado"    "Connecticut" "D.C."        "Delaware"
\end{verbatim}

\normalsize
\end{frame}

\begin{frame}[fragile]{州別の色分け (2): 地図を塗るコード}
\protect\hypertarget{ux5ddeux5225ux306eux8272ux5206ux3051-2-ux5730ux56f3ux3092ux5857ux308bux30b3ux30fcux30c9}{}
\begin{itemize}
\tightlist
\item
  作成したリストを使って、指定した州をまとめて塗りつぶします。
\end{itemize}

\scriptsize

\begin{Shaded}
\begin{Highlighting}[]
\FunctionTok{library}\NormalTok{(maps)}
\CommentTok{\# 1. まず米国本土の白地図を描く}
\FunctionTok{map}\NormalTok{(}\AttributeTok{database =} \StringTok{"state"}\NormalTok{) }

\CommentTok{\# 2. 「青い州のリスト」に含まれる州を青色(blue)で一気に塗る}
\CommentTok{\# regions 引数に名前のリストを渡す}
\FunctionTok{map}\NormalTok{(}\AttributeTok{database =} \StringTok{"state"}\NormalTok{, }\AttributeTok{regions =}\NormalTok{ blue\_states, }
    \AttributeTok{col =} \StringTok{"blue"}\NormalTok{, }\AttributeTok{fill =} \ConstantTok{TRUE}\NormalTok{, }\AttributeTok{add =} \ConstantTok{TRUE}\NormalTok{)}

\CommentTok{\# 3. 「赤い州のリスト」に含まれる州を赤色(red)で一気に塗る}
\FunctionTok{map}\NormalTok{(}\AttributeTok{database =} \StringTok{"state"}\NormalTok{, }\AttributeTok{regions =}\NormalTok{ red\_states, }
    \AttributeTok{col =} \StringTok{"red"}\NormalTok{, }\AttributeTok{fill =} \ConstantTok{TRUE}\NormalTok{, }\AttributeTok{add =} \ConstantTok{TRUE}\NormalTok{)}
\end{Highlighting}
\end{Shaded}

\normalsize
\end{frame}

\begin{frame}{州別の色分け (3): 選挙結果のマップ}
\protect\hypertarget{ux5ddeux5225ux306eux8272ux5206ux3051-3-ux9078ux6319ux7d50ux679cux306eux30deux30c3ux30d7}{}
\scriptsize

\begin{center}\includegraphics[height=0.75\textheight]{QSS_Ch05_3_files/figure-beamer/unnamed-chunk-9-1} \end{center}
\normalsize
\end{frame}

\hypertarget{ux30a6ux30a9ux30ebux30deux30fcux30c8ux306eux62e1ux5927}{%
\section{5.3.4
ウォルマートの拡大}\label{ux30a6ux30a9ux30ebux30deux30fcux30c8ux306eux62e1ux5927}}

\begin{frame}[fragile]{walmart データの読み込み (1)}
\protect\hypertarget{walmart-ux30c7ux30fcux30bfux306eux8aadux307fux8fbcux307f-1}{}
\begin{itemize}
\tightlist
\item
  拡大プロセスの視覚化のため、出店データを読み込みます。
\end{itemize}

\scriptsize

\begin{Shaded}
\begin{Highlighting}[]
\CommentTok{\# 1. ローカルに保存したデータの読み込み（推奨）}
\NormalTok{walmart }\OtherTok{\textless{}{-}} \FunctionTok{read.csv}\NormalTok{(}\StringTok{"walmart.csv"}\NormalTok{)}

\CommentTok{\# （参考）URLから直接読み込むことも可能}
\CommentTok{\# walmart \textless{}{-} read.csv("https://ayumu{-}tanaka.github.io/QSS/QSS\_Data/walmart.csv")}
\end{Highlighting}
\end{Shaded}

\normalsize
\end{frame}

\begin{frame}[fragile]{拡大プロセスの視覚化 (2): 準備}
\protect\hypertarget{ux62e1ux5927ux30d7ux30edux30bbux30b9ux306eux8996ux899aux5316-2-ux6e96ux5099}{}
\scriptsize

\begin{Shaded}
\begin{Highlighting}[]
\CommentTok{\# 1. 出店形態に応じた「色の列」をデータに追加する}
\NormalTok{walmart}\SpecialCharTok{$}\NormalTok{col }\OtherTok{\textless{}{-}} \ConstantTok{NA} \CommentTok{\# 初期化}

\CommentTok{\# 店舗タイプごとに、rgb()で透過色を指定}
\CommentTok{\# Wal{-}MartStore は赤(red)}
\NormalTok{walmart}\SpecialCharTok{$}\NormalTok{col[walmart}\SpecialCharTok{$}\NormalTok{type }\SpecialCharTok{==} \StringTok{"Wal{-}MartStore"}\NormalTok{] }\OtherTok{\textless{}{-}} \FunctionTok{rgb}\NormalTok{(}\DecValTok{1}\NormalTok{, }\DecValTok{0}\NormalTok{, }\DecValTok{0}\NormalTok{, }\FloatTok{0.2}\NormalTok{)}
\CommentTok{\# SuperCenter は緑(green)}
\NormalTok{walmart}\SpecialCharTok{$}\NormalTok{col[walmart}\SpecialCharTok{$}\NormalTok{type }\SpecialCharTok{==} \StringTok{"SuperCenter"}\NormalTok{] }\OtherTok{\textless{}{-}} \FunctionTok{rgb}\NormalTok{(}\DecValTok{0}\NormalTok{, }\DecValTok{1}\NormalTok{, }\DecValTok{0}\NormalTok{, }\FloatTok{0.2}\NormalTok{)}
\CommentTok{\# DistributionCenter は青(blue)}
\NormalTok{walmart}\SpecialCharTok{$}\NormalTok{col[walmart}\SpecialCharTok{$}\NormalTok{type }\SpecialCharTok{==} \StringTok{"DistributionCenter"}\NormalTok{] }\OtherTok{\textless{}{-}} \FunctionTok{rgb}\NormalTok{(}\DecValTok{0}\NormalTok{, }\DecValTok{0}\NormalTok{, }\DecValTok{1}\NormalTok{, }\FloatTok{0.5}\NormalTok{)}
\end{Highlighting}
\end{Shaded}

\normalsize
\end{frame}

\begin{frame}[fragile]{拡大プロセスの視覚化 (3): プロットのコード}
\protect\hypertarget{ux62e1ux5927ux30d7ux30edux30bbux30b9ux306eux8996ux899aux5316-3-ux30d7ux30edux30c3ux30c8ux306eux30b3ux30fcux30c9}{}
\scriptsize

\begin{Shaded}
\begin{Highlighting}[]
\CommentTok{\# 1. 白地図を描く}
\FunctionTok{map}\NormalTok{(}\AttributeTok{database =} \StringTok{"state"}\NormalTok{)}

\CommentTok{\# 2. 店舗の位置(long, lat)を、設定した色(col)でプロット}
\CommentTok{\# cex (点の大きさ) は、配送センター(DistributionCenter)のみ大きくする}
\FunctionTok{points}\NormalTok{(walmart}\SpecialCharTok{$}\NormalTok{long, walmart}\SpecialCharTok{$}\NormalTok{lat, }\AttributeTok{col =}\NormalTok{ walmart}\SpecialCharTok{$}\NormalTok{col, }
       \AttributeTok{pch =} \DecValTok{19}\NormalTok{, }\AttributeTok{cex =} \FunctionTok{ifelse}\NormalTok{(walmart}\SpecialCharTok{$}\NormalTok{type }\SpecialCharTok{==} \StringTok{"DistributionCenter"}\NormalTok{, }\DecValTok{1}\NormalTok{, }\FloatTok{0.5}\NormalTok{))}

\CommentTok{\# 3. 凡例を追加}
\FunctionTok{legend}\NormalTok{(}\StringTok{"bottomright"}\NormalTok{, }\AttributeTok{legend =} \FunctionTok{c}\NormalTok{(}\StringTok{"Store"}\NormalTok{, }\StringTok{"Supercenter"}\NormalTok{, }\StringTok{"Dist. Center"}\NormalTok{),}
       \AttributeTok{col =} \FunctionTok{c}\NormalTok{(}\StringTok{"red"}\NormalTok{, }\StringTok{"green"}\NormalTok{, }\StringTok{"blue"}\NormalTok{), }\AttributeTok{pch =} \DecValTok{19}\NormalTok{, }\AttributeTok{bty =} \StringTok{"n"}\NormalTok{, }\AttributeTok{cex =} \FloatTok{0.7}\NormalTok{)}
\end{Highlighting}
\end{Shaded}

\normalsize
\end{frame}

\begin{frame}{拡大プロセスの視覚化 (4): ウォルマートの拡大マップ}
\protect\hypertarget{ux62e1ux5927ux30d7ux30edux30bbux30b9ux306eux8996ux899aux5316-4-ux30a6ux30a9ux30ebux30deux30fcux30c8ux306eux62e1ux5927ux30deux30c3ux30d7}{}
\scriptsize

\begin{center}\includegraphics[height=0.75\textheight]{QSS_Ch05_3_files/figure-beamer/unnamed-chunk-13-1} \end{center}
\normalsize
\end{frame}

\hypertarget{ux307eux3068ux3081}{%
\section{5.3.5 まとめ}\label{ux307eux3068ux3081}}

\begin{frame}[fragile]{まとめ}
\protect\hypertarget{ux307eux3068ux3081-1}{}
\begin{itemize}
\tightlist
\item
  \textbf{空間分析}: 現象の地理的パターンを視覚化する強力な手法。
\item
  \textbf{maps パッケージ}: 基本的な地図描画と、特定地域の拡大が可能。
\item
  \textbf{座標とプロット}: 緯度・経度を \texttt{points()}
  等で地図上に重ねる。
\item
  \textbf{色の活用}: \texttt{rgb()} と \texttt{alpha}
  を用いて、データの密度や属性を表現する。
\item
  \textbf{コロプレスマップ}:
  地域（州など）を属性データに基づいて色分けする。
\end{itemize}
\end{frame}

\end{document}
